\chapter{Introducere}

\section{Scopul aplicației}

In contextul dezvoltarii sistemelor inteligente bazate pe viziune artificiala, extragerea și interpretarea informației vizuale din imagini reprezinta o direcție de cercetare și aplicare cu o importanța tot mai mare. Numeroase aplicații moderne depind de capacitatea sistemelor software de a identifica automat informații relevante din date vizuale captate în medii reale, unde imaginile sunt frecvent afectate de zgomot, deformari geometrice, variații de iluminare, neclaritate sau alte surse de incertitudine. în acest cadru, problema nu consta doar în detectarea unor obiecte sau simboluri, ci și în interpretarea corecta a semnificației acestora în condiții care nu pot fi controlate complet.

Scopul acestei lucrari, intitulate \textit{Sistem inteligent pentru extracția și interpretarea informației vizuale din imagini în condiții de zgomot și incertitudine}, este proiectarea și dezvoltarea unei aplicații software capabile sa prelucreze imagini dintr-un scenariu real și sa transforme informația vizuala conținută în acestea în date utile, utilizabile automat în cadrul unui sistem informatic complex. Contextul aplicativ ales este cel al managementului inteligent al accesului într-o parcare, însa contribuția lucrarii depășește acest domeniu particular, punînd accent pe metodele și mecanismele prin care informația vizuala este extrasa, reconstruita și interpretata în mod robust.

Din aceasta perspectiva, aplicația dezvoltata urmărește realizarea unui flux complet de analiza vizuala și integrare software. Sistemul include componente pentru captarea și procesarea imaginilor, identificarea regiunilor de interes, aplicarea unor algoritmi de recunoaștere, interpretarea rezultatelor și transmiterea acestora catre celelalte module ale aplicației. Daca într-o etapa anterioara accentul principal a fost pus pe recunoașterea numerelor de înmatriculare, în cadrul prezentei lucrari componenta centrala este reprezentata de detecția și decodarea codurilor QR, realizata prin combinarea metodelor clasice de procesare a imaginilor cu tehnici bazate pe învățare automata.

Scopul aplicației este, așadar, dublu. Pe de o parte, se urmărește dezvoltarea unui mecanism robust de extracție a informației vizuale din imagini degradate sau incerte, în special în cazul codurilor QR surprinse în condiții dificile de captura. Pe de alta parte, se urmărește integrarea acestui mecanism într-o arhitectura software completa, formata din backend, servicii de comunicație, baza de date, aplicație desktop și aplicație web. în acest fel, lucrarea nu trateaza exclusiv o problema algoritmica izolata, ci propune un sistem inteligent complet, în care componenta de analiza vizuala devine parte activa a unui proces operațional real.

Din punct de vedere funcțional, aplicația permite generarea și utilizarea codurilor QR asociate unor operații din sistem, detectarea lor din imagini captate sau încarcate, decodarea informației conținute și valorificarea rezultatului în logica aplicației. în paralel, sunt pastrate și extinse funcționalitași deja existente, precum procesarea imaginilor asociate numerelor de înmatriculare, comunicarea prin HTTP și WebSocket, gestionarea sigura a datelor și persistența informațiilor în baza de date. Prin urmare, scopul aplicației este de a demonstra ca informația vizuala extrasa din imagini poate fi utilizata fiabil într-un sistem software distribuit, chiar și atunci cînd datele de intrare sunt afectate de zgomot și incertitudine.

\section{Motivarea aplicației}

Motivarea acestei lucrari pornește de la o problema generala întîlnita în numeroase aplicații de viziune artificiala: imaginile provenite din scenarii reale sunt rareori ideale. în practica, informația vizuala trebuie extrasa din imagini afectate de zgomot, blur, compresie, iluminare neuniforma, reflexii, ocluzii parțiale, poziționări necontrolate ale camerei sau deformari de perspectiva. în astfel de condiții, simpla aplicare a unor metode clasice de detecție sau recunoaștere nu este întotdeauna suficienta. De aici apare necesitatea dezvoltarii unor sisteme inteligente capabile sa trateze explicit incertitudinea datelor și sa menține un nivel ridicat de robustețe în condiții de utilizare reala.

Titlul lucrarii reflecta tocmai aceasta preocupare: nu doar identificarea unor elemente vizuale, ci extracția și interpretarea informației conținute în imagini afectate de factori perturbatori. Din acest motiv, problema abordata este relevanta atît din punct de vedere teoretic, cît și practic. Din punct de vedere teoretic, ea implica îmbinarea procesarii clasice a imaginilor cu metode moderne de învățare automata. Din punct de vedere practic, ea raspunde nevoii de automatizare a unor procese reale, în care deciziile trebuie luate pe baza unor imagini imperfecte.

Alegerea contextului aplicației --- managementul accesului într-o parcare --- este motivata de faptul ca acest domeniu ofera un scenariu concret în care informația vizuala joaca un rol central. Validarea intrarii și ieșirii, asocierea unui vehicul cu o sesiune de parcare, verificarea unui tichet sau confirmarea unei operații pot fi automatizate daca sistemul este capabil sa interpreteze corect informația conținută în imagini. Totuși, într-un mediu real, camerele pot surprinde imagini cu contrast redus, mișcare, zgomot sau cadre incomplete, ceea ce transforma analiza vizuala într-o problema semnificativ mai dificila decît în condiții de laborator.

Interesul particular pentru codurile QR este justificat de rolul lor practic și de dificultatea procesarii lor în condiții neideale. Codurile QR pot reprezenta eficient tichete, identificatori de sesiune, confirmari de plata sau alte elemente esențiale în cadrul aplicației. în același timp, ele sunt sensibile la deformari de perspectiva, la degradari locale și la pierderea contrastului dintre modulele alb-negru. în situațiile în care un cod QR nu mai poate fi interpretat corect prin metode clasice, este necesara o abordare mai robusta, capabila sa reconstruiasca structura logica a codului chiar și atunci cînd imaginea bruta este afectata de incertitudine.

Aceasta observație a motivat adoptarea unei arhitecturi hibride. în locul unei soluții bazate exclusiv pe algoritmi clasici sau exclusiv pe rețele neuronale, lucrarea propune o combinație între preprocesare și rectificare geometrica, pe de o parte, și un model bazat pe învățare automata, pe de alta parte. Metodele clasice contribuie la localizarea și normalizarea regiunii de interes, în timp ce modelul neuronal intervine în etapa de reconstrucție și interpretare a matricei codului QR. Aceasta combinație este motivata de necesitatea de a trata mai eficient imaginile zgomotoase sau incerte, acolo unde o singura categorie de metode ar putea fi insuficienta.

Lucrarea este motivata și de continuitatea unui proiect anterior, dar și de necesitatea extinderii acestuia într-o direcție mai generala și mai matura. Daca soluția dezvoltata anterior era centrata în special pe recunoașterea numerelor de înmatriculare, noua versiune a sistemului pune accent pe o problema mai complexa de interpretare a informației vizuale, în care codurile QR devin elementul principal de interes. Aceasta tranziție marcheaza evoluția de la o aplicație orientata preponderent spre OCR la un sistem mai larg, capabil sa trateze diverse forme de informație vizuala și sa le integreze într-un flux software complet.

în plus, motivația aplicației este susținută și de valoarea sa practica. Un sistem capabil sa extraga și sa interpreteze automat informația vizuala din imagini afectate de zgomot și incertitudine poate reduce intervenția umana, poate accelera procesele operaționale și poate crește fiabilitatea sistemelor de acces automatizat. Integrarea acestei capabilitași într-o aplicație desktop, într-un backend de comunicație și într-o aplicație web demonstreaza ca rezultatul nu este doar un model experimental, ci o soluție software aplicabila în scenarii reale.
